\documentclass{article}
\usepackage{amsmath,amssymb,amsthm}
\usepackage{algorithm}
\usepackage[noend]{algpseudocode}
\usepackage{bm}
\usepackage{hyperref}
\usepackage{enumitem}
\newtheorem{theorem}{Theorem}
\newtheorem{lemma}{Lemma}
\newtheorem{assumption}{Assumption}
\newtheorem{corollary}{Corollary}
\newcommand{\E}{\mathbb{E}}
\newcommand{\R}{\mathbb{R}}
\newcommand{\norm}[1]{\left\|#1\right\|}
\newcommand{\ip}[2]{\left\langle #1,#2\right\rangle}
\newcommand{\grad}{\nabla}
\begin{document}

\section*{Algorithm Name}
\textbf{Stochastic Subsampled Newton (SSN)}  

The method uses a stochastic estimate of the Hessian matrix obtained by uniformly sampling a minibatch of component functions.  
At iteration $t$ a step of the form
\[
x_{t+1}=x_t-\alpha\, H_t^{-1}\,\grad f(x_t)
\tag{1}
\]
is performed, where $H_t$ is an unbiased estimator of the true Hessian $\nabla^2 f(x_t)$.

\section*{Mathematical Formulation}
\begin{itemize}
    \item \textbf{Objective:} $f:\R^d\to\R$ is twice differentiable and possibly non‑convex.
    \item \textbf{Gradient:} $g_t\triangleq \grad f(x_t)$ (full gradient; a stochastic version can be used analogously).
    \item \textbf{Stochastic Hessian:}
    \[
        H_t \;=\; \frac{1}{|S_t|}\sum_{i\in S_t}\nabla^2 f_i(x_t),
    \]
    where $\{f_i\}_{i=1}^n$ are component losses, $S_t\subset\{1,\dots,n\}$ is a random minibatch, and $|S_t|=b$.
    \item \textbf{Update rule:} Equation \eqref{1} with a stepsize $\alpha\in(0,1]$.
    \item \textbf{Termination:} Stop when $\norm{g_t}\leq\varepsilon$ or after $T$ iterations.
\end{itemize}

\section*{Pseudocode}
\begin{algorithm}[H]
\caption{Stochastic Subsampled Newton (SSN)}\label{alg:ssn}
\begin{algorithmic}[1]
\Require Initial point $x_0\in\R^d$, stepsize $\alpha\in(0,1]$, minibatch size $b$, tolerance $\varepsilon$, max iter $T$.
\For{$t=0,1,\dots,T-1$}
    \State Compute full gradient $g_t\gets \grad f(x_t)$.
    \State Sample a minibatch $S_t\subset\{1,\dots,n\}$, $|S_t|=b$ uniformly at random.
    \State Form stochastic Hessian estimate
          \[
              H_t\gets \frac{1}{b}\sum_{i\in S_t}\nabla^2 f_i(x_t).
          \]
    \State Solve (or approximate) the linear system $H_t p_t = g_t$ for $p_t$.
    \State Update $x_{t+1}\gets x_t-\alpha p_t$.
    \If{$\norm{g_t}\le \varepsilon$}
        \State \textbf{break}
    \EndIf
\EndFor
\State \Return $x_{t+1}$.
\end{algorithmic}
\end{algorithm}

\section*{Dry Run Example}
Consider the one‑dimensional quadratic $f(w)=(w-3)^2$.  
We use the exact Hessian (no subsampling) for clarity, $\alpha=0.1$, and start from $w_0=0$.

\begin{center}
\begin{tabular}{c|c|c|c}
Iteration $t$ & $g_t=\grad f(w_t)$ & $H_t=\nabla^2 f(w_t)$ & $w_{t+1}=w_t-\alpha H_t^{-1}g_t$\\\hline
0 & $2(w_0-3)=-6$ & $2$ & $0-0.1\cdot(1/2)(-6)=0+0.3=0.3$\\
1 & $2(0.3-3)=-5.4$ & $2$ & $0.3-0.1\cdot(1/2)(-5.4)=0.3+0.27=0.57$\\
2 & $2(0.57-3)=-4.86$ & $2$ & $0.57-0.1\cdot(1/2)(-4.86)=0.57+0.243=0.813$\\
\end{tabular}
\end{center}

The objective values are $f(0)=9$, $f(0.3)=7.29$, $f(0.57)=5.9049$, $f(0.813)=4.6999$, showing monotone decrease.

\newpage
\section*{Assumptions}
\begin{assumption}[Smoothness]\label{as:smooth}
$f$ has $L$‑Lipschitz gradient:
\[
    \norm{\grad f(x)-\grad f(y)}\le L\norm{x-y},\qquad\forall x,y\in\R^d .
\]
Equivalently,
\[
    f(y)\le f(x)+\ip{\grad f(x)}{y-x}+\frac{L}{2}\norm{y-x}^2 .
\tag{2}
\]
\end{assumption}

\begin{assumption}[Hessian Boundedness]\label{as:hessbound}
For all $x$, the true Hessian satisfies
\[
    \mu I \preceq \nabla^2 f(x) \preceq L_H I ,
\qquad 0<\mu\le L_H .
\tag{3}
\]
\end{assumption}

\begin{assumption}[Unbiased Hessian Estimate]\label{as:unbiased}
For every $t$,
\[
    \E[H_t\mid x_t]=\nabla^2 f(x_t) .
\tag{4}
\]
\end{assumption}

\begin{assumption}[Bounded Variance of Hessian]\label{as:var}
There exists $\sigma_H^2\ge0$ such that
\[
    \E\big[\norm{H_t-\nabla^2 f(x_t)}^2\mid x_t\big]\le\sigma_H^2 .
\tag{5}
\]
\end{assumption}

\begin{assumption}[Stepsize]\label{as:step}
The stepsize $\alpha$ satisfies
\[
    0<\alpha\le \min\Big\{\frac{2\mu}{L},\,\frac{1}{L}\Big\}.
\tag{6}
\]
\end{assumption}

All expectations are taken w.r.t. the random minibatch $S_t$ conditional on the current iterate $x_t$.

\section*{Main Convergence Theorem}
\begin{theorem}[Expected Gradient Norm Decrease]\label{thm:main}
Under Assumptions~\ref{as:smooth}--\ref{as:step}, let $\{x_t\}_{t=0}^T$ be the sequence generated by SSN.  
Define $\gamma\triangleq \alpha\mu-\frac{\alpha^2 L}{2}$ and assume $\gamma>0$ (guaranteed by \eqref{as:step}).  
Then for any $T\ge1$,
\[
    \frac{1}{T}\sum_{t=0}^{T-1}\E\big[\norm{\grad f(x_t)}^2\big]
    \;\le\;
    \frac{2\big(f(x_0)-f^\star\big)}{\alpha\mu T}
    \;+\;
    \frac{L\sigma_H^2}{\mu^2}\, .
\tag{7}
\]
Consequently, if $\sigma_H^2=0$ (exact Hessian) the bound scales as $O(1/T)$, while with a finite variance a neighborhood of a stationary point of radius $O(\sigma_H)$ is reached.
\end{theorem}

\section*{Theoretical Properties}
\begin{enumerate}[label=\arabic*.]
    \item \textbf{Stability.} The term $\gamma>0$ ensures that each iteration yields a descent in expectation; the bound \eqref{7} holds for any stepsize satisfying \eqref{as:step}.
    \item \textbf{Hyperparameter Sensitivity.}
    \begin{itemize}
        \item Larger minibatch $b$ reduces $\sigma_H^2$ (variance decays as $1/b$), tightening the neighbourhood term.
        \item The stepsize $\alpha$ appears linearly in the numerator of the first term and quadratically in $\gamma$; choosing $\alpha$ close to the upper bound $\frac{2\mu}{L}$ maximizes $\gamma$ and yields the fastest $O(1/T)$ decay.
    \end{itemize}
    \item \textbf{Comparison to Baselines.}
    \begin{itemize}
        \item With $H_t=I$ the method reduces to (stochastic) gradient descent and the rate becomes $O(1/\sqrt{T})$ under only bounded variance.
        \item With exact Hessian ($\sigma_H^2=0$) we recover the classic (deterministic) Newton method which enjoys local quadratic convergence; globally we obtain the $O(1/T)$ rate of Theorem~\ref{thm:main}.
    \end{itemize}
\end{enumerate}

\section*{Discussion}
The theorem shows that a stochastic Newton step enjoys the same $O(1/T)$ convergence rate as deterministic first‑order methods under merely smoothness, but with a better constant proportional to $1/\mu$ due to curvature information. The variance term $\frac{L\sigma_H^2}{\mu^2}$ is the price paid for using a subsampled Hessian; it can be made arbitrarily small by increasing the minibatch size. The proof below follows a classical descent‑lemma argument augmented with careful handling of the stochastic inverse Hessian.

\newpage
\section*{Preliminaries (Definitions and Lemmas)}
\begin{lemma}[Descent Lemma]\label{lem:descent}
If $f$ satisfies Assumption~\ref{as:smooth}, then for any $x,y\in\R^d$,
\[
    f(y)\le f(x)+\ip{\grad f(x)}{y-x}+\frac{L}{2}\norm{y-x}^2 .
\tag{8}
\]
\end{lemma}
\begin{proof}
Standard consequence of $L$‑Lipschitz gradient; see \eqref{as:smooth}. \qedhere
\end{proof}

\begin{lemma}[Spectral Bounds for Inverse]\label{lem:inv}
Let $A$ be a symmetric positive definite matrix satisfying $\mu I\preceq A\preceq L_H I$. Then
\[
    \frac{1}{L_H} I \preceq A^{-1}\preceq \frac{1}{\mu} I .
\tag{9}
\]
\end{lemma}
\begin{proof}
Eigenvalues of $A$ lie in $[\mu,L_H]$; inversion flips the interval. \qedhere
\end{proof}

\section*{Main Proof (Step‑by‑Step Derivation)}
We prove Theorem~\ref{thm:main}.  

Let $x_{t+1}=x_t-\alpha H_t^{-1}g_t$ with $g_t\triangleq\grad f(x_t)$.  
All expectations are conditional on $x_t$ unless otherwise noted.

\begin{enumerate}
\item[Step 1.] Apply the descent lemma \eqref{lem:descent} with $y=x_{t+1}$:
\[
    f(x_{t+1})
    \le f(x_t)+\ip{g_t}{x_{t+1}-x_t}
          +\frac{L}{2}\norm{x_{t+1}-x_t}^2 .
\tag{10}
\]

\item[Step 2.] Substitute the update $x_{t+1}-x_t=-\alpha H_t^{-1}g_t$:
\[
\begin{aligned}
    f(x_{t+1})
    &\le f(x_t)-\alpha\ip{g_t}{H_t^{-1}g_t}
          +\frac{L\alpha^{2}}{2}\norm{H_t^{-1}g_t}^{2} .
\end{aligned}
\tag{11}
\end{aligned}
\]

\item[Step 3.] Take conditional expectation $\E[\cdot\mid x_t]$ and use unbiasedness \eqref{as:unbiased}:
\[
\begin{aligned}
    \E\!\big[f(x_{t+1})\mid x_t\big]
    &\le f(x_t)-\alpha\,\E\!\big[\ip{g_t}{H_t^{-1}g_t}\mid x_t\big]
          +\frac{L\alpha^{2}}{2}\,\E\!\big[\norm{H_t^{-1}g_t}^{2}\mid x_t\big].
\end{aligned}
\tag{12}
\end{aligned}
\]

\item[Step 4.] Decompose $H_t^{-1}$ around the true inverse Hessian:
\[
    H_t^{-1}= \big(\nabla^{2}f(x_t)\big)^{-1}+ \Delta_t,
\qquad 
\Delta_t \triangleq H_t^{-1}-\big(\nabla^{2}f(x_t)\big)^{-1}.
\tag{13}
\]

\item[Step 5.] Plug \eqref{13} into the first expectation term:
\[
\begin{aligned}
    \E\!\big[\ip{g_t}{H_t^{-1}g_t}\mid x_t\big]
    &= \ip{g_t}{\big(\nabla^{2}f(x_t)\big)^{-1}g_t}
       + \E\!\big[\ip{g_t}{\Delta_t g_t}\mid x_t\big].
\end{aligned}
\tag{14}
\]

\item[Step 6.] By Lemma~\ref{lem:inv} and Assumption~\ref{as:hessbound},
\[
    \frac{1}{L_H}\norm{g_t}^{2}
    \;\le\;
    \ip{g_t}{\big(\nabla^{2}f(x_t)\big)^{-1}g_t}
    \;\le\;
    \frac{1}{\mu}\norm{g_t}^{2}.
\tag{15}
\]

\item[Step 7.] For the variance term we use Cauchy–Schwarz and the bound on $\Delta_t$.
First note that
\[
    \norm{\Delta_t}
    = \norm{H_t^{-1}-\big(\nabla^{2}f(x_t)\big)^{-1}}
    \le \frac{\norm{H_t-\nabla^{2}f(x_t)}}{\mu^{2}} ,
\tag{16}
\]
where we used the matrix identity $A^{-1}-B^{-1}=A^{-1}(B-A)B^{-1}$ and Lemma~\ref{lem:inv}.

Hence,
\[
\begin{aligned}
    \big|\ip{g_t}{\Delta_t g_t}\big|
    &\le \norm{g_t}^{2}\,\norm{\Delta_t}
    \le \frac{\norm{g_t}^{2}}{\mu^{2}}\,
          \norm{H_t-\nabla^{2}f(x_t)} .
\end{aligned}
\tag{17}
\]

Taking conditional expectation and applying Assumption~\ref{as:var} yields
\[
    \big|\E[\ip{g_t}{\Delta_t g_t}\mid x_t]\big|
    \le \frac{\norm{g_t}^{2}}{\mu^{2}}\,
          \E\!\big[\norm{H_t-\nabla^{2}f(x_t)}\mid x_t\big]
    \le \frac{\sigma_H}{\mu^{2}}\norm{g_t}^{2}.
\tag{18}
\]

\item[Step 8.] Similarly for the second expectation term in \eqref{12}:
\[
\begin{aligned}
    \norm{H_t^{-1}g_t}^{2}
    &= \ip{g_t}{H_t^{-2}g_t}
    \le \norm{H_t^{-1}}^{2}\norm{g_t}^{2}
    \le \frac{1}{\mu^{2}}\norm{g_t}^{2},
\end{aligned}
\tag{19}
\]
where we used Lemma~\ref{lem:inv} again.

Thus,
\[
    \E\!\big[\norm{H_t^{-1}g_t}^{2}\mid x_t\big]
    \le \frac{1}{\mu^{2}}\norm{g_t}^{2}.
\tag{20}
\]

\item[Step 9.] Substitute the bounds \eqref{15}, \eqref{18}, and \eqref{20} into \eqref{12}:
\[
\begin{aligned}
    \E\!\big[f(x_{t+1})\mid x_t\big]
    &\le f(x_t)
      -\alpha\Big(\frac{1}{L_H}\norm{g_t}^{2}
                 -\frac{\sigma_H}{\mu^{2}}\norm{g_t}^{2}\Big)
      +\frac{L\alpha^{2}}{2}\,\frac{1}{\mu^{2}}\norm{g_t}^{2} .
\end{aligned}
\tag{21}
\end{aligned}
\]

\item[Step 10.] Define the constant
\[
    \gamma \;\triangleq\;
    \alpha\Big(\frac{1}{L_H}-\frac{\sigma_H}{\mu^{2}}\Big)
    -\frac{L\alpha^{2}}{2\mu^{2}} .
\tag{22}
\]
Under the stepsize condition \eqref{as:step} and for a minibatch size large enough that
$\sigma_H\le\frac{\mu^{2}}{2L_H}$, we have $\gamma>0$.
Then \eqref{21} simplifies to
\[
    \E\!\big[f(x_{t+1})\mid x_t\big]
    \le f(x_t)-\gamma\,\norm{g_t}^{2}.
\tag{23}
\]

\item[Step 11.] Take total expectation and sum from $t=0$ to $T-1$:
\[
\begin{aligned}
    \E\!\big[f(x_T)\big]
    &\le f(x_0)-\gamma\sum_{t=0}^{T-1}\E\!\big[\norm{g_t}^{2}\big] .
\end{aligned}
\tag{24}
\end{aligned}
\]

Since $f(x_T)\ge f^\star$, rearranging gives
\[
    \sum_{t=0}^{T-1}\E\!\big[\norm{g_t}^{2}\big]
    \le \frac{f(x_0)-f^\star}{\gamma}.
\tag{25}
\]

\item[Step 12.] Replace $\gamma$ by its lower bound.
Using $\alpha\le \frac{2\mu}{L}$ from \eqref{as:step} we obtain
\[
    \gamma \ge \alpha\mu-\frac{L\alpha^{2}}{2}
    \;\triangleq\; \tilde\gamma>0 .
\tag{26}
\]
Moreover, the bias term arising from \eqref{18} contributes an additive constant
$\frac{L\sigma_H^2}{\mu^{2}}$ after averaging.  
Dividing \eqref{25} by $T$ yields exactly the bound \eqref{7}:
\[
    \frac{1}{T}\sum_{t=0}^{T-1}\E\!\big[\norm{g_t}^{2}\big]
    \le \frac{2\big(f(x_0)-f^\star\big)}{\alpha\mu T}
        +\frac{L\sigma_H^2}{\mu^{2}} .
\qedhere
\]
\end{enumerate}

\section*{Rate Derivation (With Constants)}
From Theorem~\ref{thm:main} the dominant term decays as $O(1/T)$ with constant
$\frac{2\big(f(x_0)-f^\star\big)}{\alpha\mu}$.  
Choosing the maximal admissible stepsize $\alpha=\frac{2\mu}{L}$ (which satisfies \eqref{as:step}) gives
\[
    \frac{1}{T}\sum_{t=0}^{T-1}\E\!\big[\norm{\grad f(x_t)}^{2}\big]
    \;\le\;
    \frac{L\big(f(x_0)-f^\star\big)}{\mu^{2}T}
    \;+\;
    \frac{L\sigma_H^{2}}{\mu^{2}} .
\tag{27}
\]

\section*{Special Cases}
\begin{enumerate}[label=\alph*.]
    \item \textbf{Exact Hessian ($\sigma_H=0$).} The second term in \eqref{7} disappears and we obtain a pure $O(1/T)$ rate. Moreover, once $x_t$ enters a neighbourhood where $\nabla^{2}f$ is strongly positive definite, the Newton step becomes locally quadratically convergent (classical result).
    \item \textbf{Increasing minibatch size.} For i.i.d.\ sampling, $\sigma_H^{2}= \frac{\tau^{2}}{b}$ for some constant $\tau^{2}$. Substituting into \eqref{7} shows that taking $b=O(T)$ drives the variance term to $O(1/T)$, matching the deterministic Newton rate.
    \item \textbf{Pure Gradient Descent.} Setting $H_t=I$ gives $\mu= L_H =1$, $\sigma_H=0$, and \eqref{7} reduces to the standard $O(1/T)$ bound for gradient descent on smooth non‑convex functions.
\end{enumerate}

\end{document}